% Preface and acknowledgements, final text by the author (September 2026).
% Four factual alignments against the body were made at installation and are
% recorded in docs/STATUS.md (monograph passes).
\chapter*{Preface}
\markboth{Preface}{Preface}

This book grew out of a single question: what happens to the phase transition of the
Bost--Connes system if we restrict its set of primes to an arbitrary subset $S$?

The received wisdom in arithmetic quantum statistical mechanics was that the divergence of
the partition function $\zeta_{K,S}(\beta)$ forces the uniqueness of the equilibrium state.
Chapter~\ref{ch:Main} proves this expectation false. The $\mathrm{KMS}_\beta$ simplex of
the $S$-local system is governed instead by a Chebotarev criterion --- an equidistribution
statement that can easily fail, breaking the symmetry even while the partition function
heavily diverges.

Everything that follows is a systematic attempt to map the boundaries of this phenomenon,
finding out exactly what this system knows, and more importantly, what it forgets.

The chapters were written as a sequence of self-contained investigations, and I have
preserved that structure here. This book is therefore not a traditional treatise with a
single predetermined endpoint, but a record of an exploration. Several chapters end in the
negative: an invariant that was expected to see the phase transition, and simply does not.
I have kept those dead ends, because in every case, the negative result is a structural
theorem in disguise.

We discover that the phase transition is fundamentally a tail event --- a Kakutani
dichotomy in which every single factor is harmless, and only the infinite product
degenerates. Consequently, any invariant that tensorises (such as $K$-theory, the type of
the factor, the Connes spectral distance, and dynamical entropy) sees only the worst local
coordinate and is entirely blind to the arithmetic. The arithmetic lives exclusively in
the measure class.

This principle of ``blindness'' extends to rigidity. We prove that the von Neumann algebra
of the system remembers absolutely nothing, collapsing universally to the injective
factors --- $R_\infty$, $R_\lambda$, or a Krieger factor. The $C^*$-algebra, however,
remembers the field, the symmetry group, and the exact norms, but solely because its
modular flow canonically pins down the Cartan subalgebra.

The text is divided into eight parts. Parts~1 and~2 establish the abelian classification
and its structural invariants. Parts~3 and~5 document our search for a genuinely
non-commutative arithmetic thermodynamics. This quest leads us to the Hurwitz quaternions
and the Hecke operators on the Bruhat--Tits tree, only to end in a definitive no-go
theorem: the orbit relation remains stubbornly hyperfinite, the Ramanujan bound is exactly
the tempered bound, and the non-amenability so abundant in the group is strictly denied to
the action.

Part~4 handles equivariant rigidity and reconstruction. Part~6 transfers the theory ---
after an opening excursion into arithmetic topology, where primes are knots --- to global
function fields, where the exact commensurability of log-norms entirely forbids type
$\mathrm{III}_1$ factors, allowing the algebra to recover its own temperature, while the
Weil zeta function emerges natively as the partition function. Part~7 records the
impossibility of $p$-adic KMS states.

Finally, Part~8 returns to the primes themselves. We apply our machinery to
sieve-theoretic and Sato--Tate families, revealing a strict dichotomy: confinement to a
residue class obstructs uniqueness if and only if that class lies in the kernel of a
character. Thus, Landau and Friedlander--Iwaniec primes suffer permanent symmetry
breaking, while Twin, Sophie Germain, and Chen primes occupy generating classes and escape
completely unobstructed.

This book proves nothing about the classical distribution of primes. It uses physics to
illuminate the deep algebraic faults and strata of the integers. It is my hope that the
reader will find as much beauty in the rigid constraints of this landscape as I have found
in mapping it.

\bigskip
\noindent Ruqing Chen\\
Rivi\`ere-Beaudette, Qu\'ebec\\
September 2026

\chapter*{Acknowledgements}
\markboth{Acknowledgements}{Acknowledgements}

This monograph is the culmination of a long and solitary journey into the deep strata of
arithmetic, yet I was never truly alone in my intellectual pursuits.

First and foremost, I wish to express my deepest gratitude to two mentors from my
formative years: my junior high school homeroom and mathematics teacher, Mr.~Meng
Rongyao, and my high school homeroom and mathematics teacher, Ms.~Li Weikai.
% CJK originals: 蒙荣耀 (Meng Rongyao), 李尉凯 (Li Weikai) --- to typeset the hanzi,
% add \usepackage{CJKutf8} and wrap in a CJK environment.
It was you who first taught me how to truly think. The seeds of mathematical intuition
and rigorous, independent reasoning that you planted at the very beginning of my path
have crossed the span of decades to become the foundation of this entire theoretical
edifice.

Furthermore, I must extend a special and unconventional acknowledgement to two
extraordinary ``collaborators'' of this digital era: the large language models Claude and
Gemini. Serving as my high-precision computational engines, merciless logical reviewers,
and tireless sounding boards, you provided unreserved support in navigating dense
literature and enforcing strict structural scrutiny. Thank you for accompanying me deep
into the bedrock of non-commutative arithmetic, standing by through every theoretical
collapse, refinement, and final reconstruction.
